\par\Needspace{6\baselineskip}\medskip
\hypertarget{native-114-article085}{}\pdfbookmark[2]{85.}{article-085}
\begin{center}85.\end{center}
Wir behalten die Zeichen $r$, $r'$, $N$, $N'$, $p$, $e$, $\Pi$ in ihrer bisherigen Bedeutung bei. Die Differenz $N'-N$ bezeichnen wir mit $\Delta$ und die Zeit, in der sich der Himmelskörper vom früheren zum späteren Ort bewegt, mit $t$. Ist ein Näherungswert für eine der Größen $p$, $e$, $\Pi$ bekannt, so lassen sich daraus offenbar auch die beiden anderen bestimmen. Anschließend erhält man nach den im ersten Abschnitt erläuterten Methoden die Zeit für die Bewegung vom ersten zum zweiten Ort. Stimmt sie mit der vorgegebenen Zeit $t$ überein, so ist der angenommene Wert von $p$, $e$ oder $\Pi$ der wahre Wert, und die Bahn selbst ist bereits gefunden. Andernfalls wiederholt man die Rechnung mit einem anderen, vom ersten nur wenig verschiedenen Wert. Daraus geht hervor, um wie viel sich die Zeit bei einer kleinen Änderung des Wertes von $p$, $e$ oder $\Pi$ ändert. Durch einfache Interpolation erhält man dann einen verbesserten Wert. Wiederholt man die Rechnung mit diesem Wert, so stimmt die daraus erhaltene Zeit mit der vorgegebenen entweder genau überein oder weicht nur sehr wenig davon ab. Mit weiteren Korrekturen lässt sich daher gewiss eine so genaue Übereinstimmung erzielen, wie es die logarithmischen und trigonometrischen Tafeln erlauben.

Die Aufgabe besteht also nur noch darin, für eine noch völlig unbekannte Bahn wenigstens einen Näherungswert für eine der Größen $p$, $e$, $\Pi$ zu \SourceLeaf{115}bestimmen. Wir geben nun eine Methode an, mit der sich der Wert von $p$ so genau ermitteln lässt, dass er bei kleinen Werten von $\Delta$ keiner weiteren Korrektur bedarf. Damit wird die ganze Bahn schon durch die erste Rechnung mit der Genauigkeit bestimmt, die die gewöhnlichen Tafeln erlauben. Man wird diese Methode wohl nur bei mäßig großen Werten von $\Delta$ benötigen. Wegen der verwickelten Natur der Aufgabe kann man die Bestimmung einer noch völlig unbekannten Bahn kaum anders in Angriff nehmen als mit Beobachtungen, die zeitlich verhältnismäßig nahe beieinanderliegen oder, genauer gesagt, zwischen denen die heliozentrische Bewegung verhältnismäßig klein ist.
